\chapter{Girsanov and Changes of Numeraire}\label{ch:qm:girsanov-and-changes-of-numeraire}

A risk system prices a five-year caplet by simulating the short rate week by week, integrating
it into a discount factor, and averaging over paths: 260 random numbers per path and a
standard error of 0.05 basis point after 20\,000 paths. A colleague prices the same caplet by
changing the measure to the one attached to the zero-coupon bond maturing at the fixing date:
under it the short rate at that date is a single Gaussian with a known mean, so each path costs
one random number, and the same argument yields a closed form that takes a microsecond. The
answers agree within their standard errors: 3.26 basis points of notional. Nothing about the
market changed; only the unit in which values were counted. This chapter proves Girsanov's
theorem, which says what a \omterm{def:qm:probability-at-speed:change}{change of measure} does to a \omterm{def:qm:brownian-motion:bm}{Brownian motion}, defines \omterm{def:qm:girsanov-and-changes-of-numeraire:numeraire}{numeraires} and
the \omterm{def:qm:probability-at-speed:martingale}{martingale} measures attached to them, and derives the forward and \omterm{def:qm:girsanov-and-changes-of-numeraire:annuity}{annuity measures} that
every rates and options chapter of the series computes with.

\section{Girsanov's theorem}

\begin{definition}[Novikov's condition]\label{def:qm:girsanov-and-changes-of-numeraire:novikov}
An \omterm{def:qm:probability-at-speed:filtration}{adapted process} $\gamma$ satisfies \emph{Novikov's condition}\index{Novikov's condition} on $[0,
T]$ if $\E[\exp(\tfrac12\int_0^T\gamma_s^2\,ds)] < \infty$.
\end{definition}

It is a sufficient condition for the positive \omterm{def:qm:ito-calculus:local}{local martingale} $Z_t = \mathcal E(\int\gamma\,dW)_t =
\exp(\int_0^t\gamma_s\,dW_s - \tfrac12\int_0^t\gamma_s^2\,ds)$ of \cref{def:qm:ito-calculus:exp} to be a true
\omterm{def:qm:probability-at-speed:martingale}{martingale}, with $\E[Z_T] = 1$, so that it can serve as a \omterm{def:qm:probability-at-speed:change}{density process}
(\cref{def:qm:probability-at-speed:change}). Bounded $\gamma$ always qualifies.

\begin{theorem}[Girsanov]\label{thm:qm:girsanov-and-changes-of-numeraire:girsanov}
Let $\gamma$ satisfy \omterm{def:qm:girsanov-and-changes-of-numeraire:novikov}{Novikov's condition} and define $\mathbb Q$ on $\mathcal F_T$ by $d\mathbb Q/d\P = Z_T$. Then
\[
W^{\mathbb Q}_t = W_t - \int_0^t\gamma_s\,ds
\]
is a \omterm{def:qm:brownian-motion:bm}{Brownian motion} under $\mathbb Q$. Conversely, every $\mathbb Q \sim \P$ on the \omterm{def:qm:probability-at-speed:filtration}{filtration} of $W$ has a density
of this form.
\end{theorem}

\begin{proof}
By \cref{thm:qm:ito-calculus:levy} it suffices that $W^{\mathbb Q}$ be a continuous $\mathbb Q$-local \omterm{def:qm:probability-at-speed:martingale}{martingale} with
$[W^{\mathbb Q}]_t = t$; the \omterm{def:qm:brownian-motion:qv}{quadratic variation} is that of $W$, since the drift has finite variation.
By \cref{prop:qm:probability-at-speed:bayes}, $W^{\mathbb Q}$ is a $\mathbb Q$-local \omterm{def:qm:probability-at-speed:martingale}{martingale} if $ZW^{\mathbb Q}$ is a
$\P$-local \omterm{def:qm:probability-at-speed:martingale}{martingale}, and by the product rule, with $dZ = Z\gamma\,dW$,
$d(ZW^{\mathbb Q}) = Z\,dW^{\mathbb Q} + W^{\mathbb Q}\,dZ + d[Z, W^{\mathbb Q}] = Z\,dW - Z\gamma\,dt + W^{\mathbb Q}\,dZ + Z\gamma\,dt$, a
stochastic integral. The converse is the \omterm{def:qm:probability-at-speed:martingale}{martingale} representation theorem applied to the
\omterm{def:qm:probability-at-speed:change}{density process} (\cref{thm:qm:ito-calculus:representation}).
\end{proof}

A \omterm{def:qm:probability-at-speed:change}{change of measure} cannot change volatility, which is visible in the \omterm{def:qm:brownian-motion:qv}{quadratic variation} of a
single path; it changes drift, which is not. With constant $\gamma$ the theorem is the
Cameron--Martin formula: reweighting paths by $e^{\gamma W_T - \gamma^2T/2}$ turns a \omterm{def:qm:brownian-motion:bm}{Brownian motion} into one
with drift $\gamma$ (\cref{fig:qm:girsanov-and-changes-of-numeraire:girsanov}). Read backwards, it removes a
drift, which is how the first-passage probability with drift of chapter~2 is derived
(\cref{exo:qm:girsanov-and-changes-of-numeraire:7}).

\begin{omfigure}
\begin{tikzpicture}
\begin{axis}[width=11cm,height=4.8cm,
  xlabel={$x$},ylabel={density of $W_1$},
  tick label style={font=\small},label style={font=\small},
  xmin=-4,xmax=5,ymin=0,ymax=0.45,grid=major,grid style={gray!25},
  legend columns=2,legend style={font=\footnotesize,at={(0.5,-0.3)},anchor=north,draw=none,fill=none,/tikz/every even column/.append style={column sep=8pt}},
  table/col sep=comma]
\addplot[omDef,thick,const plot mark mid,mark=none] table[x=x,y=p]{figdata/methods/05-girsanov-and-changes-of-numeraire/girsanov.csv};
\addplot[omThm,thick,const plot mark mid] table[x=x,y=q]{figdata/methods/05-girsanov-and-changes-of-numeraire/girsanov.csv};
\addplot[omDef,dashed,domain=-4:5,samples=80,forget plot] {exp(-x^2/2)/sqrt(2*pi)};
\addplot[omThm,dashed,domain=-4:5,samples=80,forget plot] {exp(-(x-1)^2/2)/sqrt(2*pi)};
\legend{{under $\P$: $\mathcal N(0,1)$},{reweighted by $e^{W_1 - 1/2}$: $\mathcal N(1,1)$}}
\end{axis}
\end{tikzpicture}

\omcaption{The same 200\,000 draws of $W_1$, counted plainly (blue) and weighted by the density
$Z_1 = e^{\gamma W_1 - \gamma^2/2}$ with $\gamma = 1$ (red): the weighted histogram is the $\mathcal N(1, 1)$ density
(dashed), and the shape is untouched. Data: the chapter's tutorial, seeded.}\label{fig:qm:girsanov-and-changes-of-numeraire:girsanov}
\end{omfigure}

\section{Numeraires and their martingale measures}

\begin{definition}[Numeraire, money-market account]\label{def:qm:girsanov-and-changes-of-numeraire:numeraire}
A \emph{numeraire}\index{numeraire} $\mathcal N$ is the price process of a traded asset (or self-financing
portfolio) that is strictly positive at all times; values expressed in units of it are
$V_t/\mathcal N_t$. The \emph{money-market account}\index{money-market account} is the numeraire $B_t =
\exp\int_0^tr_s\,ds$ that rolls overnight at the short rate $r_t$.
\end{definition}

\begin{definition}[Equivalent martingale measure, risk-neutral measure]\label{def:qm:girsanov-and-changes-of-numeraire:emm}
An \emph{equivalent martingale measure}\index{equivalent martingale measure} for a \omterm{def:qm:girsanov-and-changes-of-numeraire:numeraire}{numeraire} $\mathcal N$
is a probability $\mathbb Q^{\mathcal N}$ equivalent to $\P$ under which $V_t/\mathcal N_t$ is a \omterm{def:qm:probability-at-speed:martingale}{martingale} for
the price $V$ of every traded asset. The \emph{risk-neutral measure}\index{risk-neutral measure}
$\mathbb Q$ is the equivalent \omterm{def:qm:probability-at-speed:martingale}{martingale} measure for the \omterm{def:qm:girsanov-and-changes-of-numeraire:numeraire}{money-market account}.
\end{definition}

Under such a measure a claim paying $V_T$ at $T$ is worth
\[
V_t = \mathcal N_t\,\E^{\mathbb Q^{\mathcal N}}_t\Bigl[\frac{V_T}{\mathcal N_T}\Bigr], \qquad\text{and with } \mathcal N = B:\quad V_t =
\E^{\mathbb Q}_t\bigl[e^{-\int_t^Tr_s\,ds}V_T\bigr].
\]
Why such a measure should exist is an economic statement: in a market without arbitrage there
is at least one, and in a complete market exactly one. That is the fundamental theorem of asset
pricing, and it is One Quant Book~5, chapter~1, that proves it; here the measures are taken as
given and used as tools. Under $\mathbb Q$, a stock that pays no dividend has drift $r_t$: $dS/S = r_t\,dt +
\sigma\,dW^{\mathbb Q}$, whatever its expected return in the world, because Girsanov with $\gamma = -(\mu -
r)/\sigma$ removes the excess.

\section{Changing the numeraire}

\begin{definition}[Change of numeraire]\label{def:qm:girsanov-and-changes-of-numeraire:changenum}
A \emph{change of numeraire}\index{change of numeraire} replaces one \omterm{def:qm:girsanov-and-changes-of-numeraire:numeraire}{numeraire} and its \omterm{def:qm:girsanov-and-changes-of-numeraire:emm}{equivalent martingale measure} by another \omterm{def:qm:girsanov-and-changes-of-numeraire:numeraire}{numeraire} and the measure under which prices divided by it are
\omterm{def:qm:probability-at-speed:martingale}{martingales}; prices are unchanged, only the unit and the measure in which they are computed change.
\end{definition}

\begin{theorem}[Change of numeraire]\label{thm:qm:girsanov-and-changes-of-numeraire:change}
If $\mathbb Q^{\mathcal N}$ is an \omterm{def:qm:girsanov-and-changes-of-numeraire:emm}{equivalent martingale measure} for $\mathcal N$ and $\mathcal M$ is another
\omterm{def:qm:girsanov-and-changes-of-numeraire:numeraire}{numeraire}, the measure defined by
\[
\frac{d\mathbb Q^{\mathcal M}}{d\mathbb Q^{\mathcal N}}\bigg|_{\mathcal F_t} = \frac{\mathcal M_t/\mathcal M_0}{\mathcal N_t/\mathcal N_0}
\]
is an \omterm{def:qm:girsanov-and-changes-of-numeraire:emm}{equivalent martingale measure} for $\mathcal M$. If $d\mathcal M/\mathcal M$ and $d\mathcal N/\mathcal N$ have
volatility vectors $\sigma_{\mathcal M}$, $\sigma_{\mathcal N}$ against a \omterm{def:qm:brownian-motion:bm}{Brownian motion} $W^{\mathcal N}$ of $\mathbb
Q^{\mathcal N}$, then $dW^{\mathcal M} = dW^{\mathcal N} - (\sigma_{\mathcal M} - \sigma_{\mathcal N})\,dt$.
\end{theorem}

\begin{proof}
The \omterm{def:qm:probability-at-speed:change}{density process} $Z_t = (\mathcal M_t/\mathcal N_t)(\mathcal N_0/\mathcal M_0)$ is a positive $\mathbb Q^{\mathcal N}$-martingale
because $\mathcal M$ is traded. For a traded $V$, $Z_t\cdot V_t/\mathcal M_t = (\mathcal N_0/\mathcal M_0)\,V_t/\mathcal N_t$ is a
$\mathbb Q^{\mathcal N}$-martingale, so $V/\mathcal M$ is a $\mathbb Q^{\mathcal M}$-martingale by
\cref{prop:qm:probability-at-speed:bayes}. By Itô's formula the ratio $\mathcal M/\mathcal N$ has volatility
$\sigma_{\mathcal M} - \sigma_{\mathcal N}$, so $Z = \mathcal E(\int(\sigma_{\mathcal M} - \sigma_{\mathcal N})\,dW^{\mathcal N})$ and Girsanov
gives the drift.
\end{proof}

Geman, El Karoui and Rochet (1995) made this a working tool. The art is to choose the
\omterm{def:qm:girsanov-and-changes-of-numeraire:numeraire}{numeraire} so that the ratio $V_T/\mathcal N_T$ is simple. \Cref{fig:qm:girsanov-and-changes-of-numeraire:map}
lists the choices the series uses; Margrabe's option to exchange one asset for another, priced
with the second asset as \omterm{def:qm:girsanov-and-changes-of-numeraire:numeraire}{numeraire}, is \cref{exo:qm:girsanov-and-changes-of-numeraire:6}.

\begin{omfigure}
\begin{tikzpicture}[font=\small]
\matrix[matrix of nodes,nodes={align=left,anchor=west,inner sep=3pt},
  column 1/.style={nodes={text width=3.6cm}},column 2/.style={nodes={text width=3.4cm}},
  column 3/.style={nodes={text width=5.6cm}},row sep=4pt,column sep=6pt] (m) {
\textbf{numeraire} & \textbf{measure} & \textbf{martingales (prices$\,/\,$numeraire)} \\
money-market account $B_t$ & risk-neutral $\mathbb Q$ & every discounted price $V_t/B_t$ \\
zero-coupon bond $P(t,T)$ & $T$-forward $\mathbb Q^T$ & forward prices $V_t/P(t,T)$; a simple rate $F(t;T{-}\delta,T)$ paid at $T$ \\
annuity $A_t$ & annuity $\mathbb Q^A$ & the par swap rate $S_t$ of the swap with fixed dates $T_i$ \\
a stock $S_t$ & share measure $\mathbb Q^S$ & $V_t/S_t$: the exchange option's ratio $S^2_t/S^1_t$ \\
};
\draw[omDef,thick] (m-1-1.north west) -- (m-1-3.north east);
\draw[omDef] (m-1-1.south west) -- (m-1-3.south east);
\draw[omDef,thick] (m-5-1.south west) -- (m-5-3.south east);
\end{tikzpicture}

\omcaption{\omterm{def:qm:girsanov-and-changes-of-numeraire:numeraire}{Numeraires} the series uses, their \omterm{def:qm:girsanov-and-changes-of-numeraire:emm}{equivalent martingale measures}, and what becomes a
\omterm{def:qm:probability-at-speed:martingale}{martingale} under each: pricing a claim means choosing the \omterm{def:qm:girsanov-and-changes-of-numeraire:numeraire}{numeraire} that makes the claim's ratio
simplest.}\label{fig:qm:girsanov-and-changes-of-numeraire:map}
\end{omfigure}

\section{The forward measure}

\begin{definition}[Forward measure]\label{def:qm:girsanov-and-changes-of-numeraire:forward}
The \emph{forward measure}\index{forward measure} $\mathbb Q^T$ for maturity $T$ is the \omterm{def:qm:girsanov-and-changes-of-numeraire:emm}{equivalent martingale
measure} for the \omterm{def:qm:girsanov-and-changes-of-numeraire:numeraire}{numeraire} $P(t, T)$, the zero-coupon bond maturing at $T$.
\end{definition}

\begin{proposition}[Discounting inside the expectation, and out]\label{prop:qm:girsanov-and-changes-of-numeraire:forward}
For a payoff $X$ at $T$: $\E^{\mathbb Q}_t[e^{-\int_t^Tr_s\,ds}X] = P(t, T)\,\E^{T}_t[X]$. The forward price
$V_t/P(t,T)$ of any traded asset is a $\mathbb Q^T$-martingale, and in particular $\E^T[r_T] = f(0, T)$,
the instantaneous forward rate.
\end{proposition}

\begin{proof}
Apply the pricing formula with $\mathcal N = P(\cdot, T)$, for which $P(T, T) = 1$. For the last claim,
$f(0,T) = -\partial_T\ln P(0,T) = \E^{\mathbb Q}[r_Te^{-\int_0^Tr}]/P(0,T) = \E^T[r_T]$.
\end{proof}

The expectation of a stochastic discount factor times a payoff has become a deterministic
discount factor times an expectation. For the Ornstein--Uhlenbeck short rate of chapter~4
($r_0 = 3\%$, $\bar r = 4\%$, $\kappa = 0.5$, $\sigma = 1\%$), the bond has volatility $-\sigma B(T - t)$, so by
\cref{thm:qm:girsanov-and-changes-of-numeraire:change} the short rate's drift under $\mathbb Q^T$ loses
$\sigma^2B(T - t)$: $r_T$ remains Gaussian, with a lower mean. $\E^{\mathbb Q}[r_5] = 3.918\%$ while $\E^5[r_5] = f(0,
5) = 3.901\%$. The gap, a convexity effect, grows to $\sigma^2/2\kappa^2 = 2$ basis points at long maturities
(\cref{fig:qm:girsanov-and-changes-of-numeraire:convexity}).

\begin{omfigure}
\begin{tikzpicture}
\begin{axis}[width=11cm,height=4.6cm,
  xlabel={maturity $t$ (years)},ylabel={$\E^{\mathbb Q}[r_t] - f(0,t)$ (bp)},
  tick label style={font=\small},label style={font=\small},
  xmin=0,xmax=30,ymin=0,ymax=2.2,grid=major,grid style={gray!25},
  table/col sep=comma]
\addplot[omDef,very thick] table[x=t,y=gap]{figdata/methods/05-girsanov-and-changes-of-numeraire/convexity.csv};
\addplot[gray,dashed,domain=0:30] {2};
\node[font=\footnotesize,anchor=north east,fill=white,inner sep=1pt] at (axis cs:29,1.9) {$\sigma^2/2\kappa^2 = 2$ bp};
\end{axis}
\end{tikzpicture}

\omcaption{Expected short rate under the \omterm{def:qm:girsanov-and-changes-of-numeraire:emm}{risk-neutral measure} minus the instantaneous forward rate,
which is the expected short rate under the \omterm{def:qm:girsanov-and-changes-of-numeraire:forward}{forward measure} of the same maturity, for the
chapter's Ornstein--Uhlenbeck short rate. The gap, $\sigma^2B(t)^2/2$, is the price of the bond's
convexity. Data: closed forms, the chapter's tutorial.}\label{fig:qm:girsanov-and-changes-of-numeraire:convexity}
\end{omfigure}

The caplet of the hook pays $\delta(F - K)^+$ at $T_1 + \delta$ on the rate $F$ fixed at $T_1 = 5$ for $\delta = \tfrac14$
year, with $K = 4.5\%$. Its value at $T_1$ is $(1 + \delta K)(x - P(T_1, T_1 + \delta))^+$ with $x = 1/(1 + \delta K) =
0.98888$: a put on a zero-coupon bond. Under $\mathbb Q^{T_1}$, $P(T_1, T_1 + \delta) = e^{A - Br_{T_1}}$ is lognormal,
and a Black-type formula follows (Jamshidian, 1989) with bond volatility $\sigma_P =
\sigma\sqrt{(1 - e^{-2\kappa T_1})/2\kappa}\,B(\delta) = 0.234\%$: 3.26 basis points of notional, USD~32\,600 on
USD~100 million. \Cref{fig:qm:girsanov-and-changes-of-numeraire:mc} compares the two simulations.

\begin{omfigure}
\begin{tikzpicture}
\begin{axis}[width=11cm,height=5cm,xmode=log,log basis x=2,
  xlabel={paths},ylabel={caplet (bp of notional)},
  tick label style={font=\small},label style={font=\small},
  xmin=180,xmax=90000,ymin=1.8,ymax=4.4,grid=major,grid style={gray!25},
  legend columns=3,legend style={font=\footnotesize,at={(0.5,-0.3)},anchor=north,draw=none,fill=none,/tikz/every even column/.append style={column sep=8pt}},
  table/col sep=comma]
\addplot[omDef,thick,mark=*,mark size=1.5pt,error bars/.cd,y dir=both,y explicit]
  table[x=n,y=q,y error plus expr=\thisrow{qhi}-\thisrow{q},y error minus expr=\thisrow{q}-\thisrow{qlo}]{figdata/methods/05-girsanov-and-changes-of-numeraire/convergence.csv};
\addplot[omThm,thick,mark=triangle*,mark size=1.8pt,error bars/.cd,y dir=both,y explicit]
  table[x expr=\thisrow{n}*1.12,y=fwd,y error plus expr=\thisrow{fhi}-\thisrow{fwd},y error minus expr=\thisrow{fwd}-\thisrow{flo}]{figdata/methods/05-girsanov-and-changes-of-numeraire/convergence.csv};
\addplot[black,dashed,domain=180:90000] {3.2586};
\legend{under $\mathbb Q$ (260 normals a path),under $\mathbb Q^{T_1}$ (one normal),closed form}
\end{axis}
\end{tikzpicture}

\omcaption{The caplet by Monte Carlo under the \omterm{def:qm:girsanov-and-changes-of-numeraire:emm}{risk-neutral measure} (weekly steps, money-market
discounting) and under the fixing-date \omterm{def:qm:girsanov-and-changes-of-numeraire:forward}{forward measure}, with $\pm 2$ standard errors, against the
closed form 3.26 bp. At equal numbers of paths the errors are the same; the forward-measure path costs
one normal instead of 260. Data: the chapter's tutorial, seeded.}\label{fig:qm:girsanov-and-changes-of-numeraire:mc}
\end{omfigure}

The gain is not a smaller variance per path: the discount factor varies little next to the
payoff, and the two standard errors per path are within half a percent of each other. It is the
dimension of the problem, 260 normals down to one, and then no simulation at all.

\section{The annuity measure}

\begin{definition}[Annuity measure]\label{def:qm:girsanov-and-changes-of-numeraire:annuity}
For a swap with fixed payments at $T_1 < \dots < T_n$ and accrual fractions $\delta_i$, the \emph{annuity
measure}\index{annuity measure} $\mathbb Q^A$ is the \omterm{def:qm:girsanov-and-changes-of-numeraire:emm}{equivalent martingale measure} for the \omterm{def:qm:girsanov-and-changes-of-numeraire:numeraire}{numeraire} $A_t =
\sum_i\delta_iP(t, T_i)$, the annuity of One Quant Book~2, chapter~9.
\end{definition}

The par swap rate is $S_t = (P(t, T_0) - P(t, T_n))/A_t$, a traded value divided by the \omterm{def:qm:girsanov-and-changes-of-numeraire:numeraire}{numeraire}:
a $\mathbb Q^A$-martingale. A payer swaption exercised at $T_0$ pays $A_{T_0}(S_{T_0} - K)^+$, so its value is $A_0\,\E^A[(S_{T_0}
- K)^+]$: with a normal law for $S_{T_0}$ this is the Bachelier formula of One Quant Book~2,
chapter~13, with a lognormal one Black's. The \omterm{def:qm:girsanov-and-changes-of-numeraire:changenum}{change of numeraire} replaces a swaption, a claim on
a whole curve, by a call on one \omterm{def:qm:probability-at-speed:martingale}{martingale}. One Quant Book~6 builds its rates models on this
measure and on the \omterm{def:qm:girsanov-and-changes-of-numeraire:forward}{forward measure}.

\section{Tutorial: two measures, one price}\label{tut:qm:girsanov-and-changes-of-numeraire:caplet}

\begin{tutorial}
\textbf{Goal.} Price the caplet under two measures, check the reweighting between them, and see
Girsanov at work on a histogram. \textbf{End state:}
\cref{fig:qm:girsanov-and-changes-of-numeraire:girsanov,fig:qm:girsanov-and-changes-of-numeraire:convexity,fig:qm:girsanov-and-changes-of-numeraire:mc}
and the price 3.26 bp.
\begin{tutsteps}
\item \textbf{The weights.} The running project converts expectations between \omterm{def:qm:girsanov-and-changes-of-numeraire:numeraire}{numeraires} and builds
Girsanov densities.
\omcode{code/firm/numeraire/firm_numeraire.py}{21}{37}{Change-of-numeraire weights, a weighted mean, and a Girsanov density.}\label{lst:qm:girsanov-and-changes-of-numeraire:weights}
\item \textbf{The forward-measure sampler}: the mean of $r_{T}$ loses $\sigma^2\int_0^Te^{-\kappa(T-s)}B(T - s)\,ds$.
\omcode{code/methods/05-girsanov-and-changes-of-numeraire/python/qm_numeraire.py}{85}{99}{The short rate at the fixing date under the forward measure, and the caplet by one normal a path.}\label{lst:qm:girsanov-and-changes-of-numeraire:forward}
\item \textbf{Run} \texttt{problem()}, \texttt{reweighting\_check()} (100\,000 risk-neutral paths
reweighted by $d\mathbb Q^{T}/d\mathbb Q$ give $\E^{T}[r_5] = 3.906\% \pm 0.003\%$ against $3.901\%$) and
\texttt{fig\_numeraire.py}.
\end{tutsteps}
\textbf{What to change next.} Price the caplet under the payment-date measure $\mathbb Q^{T_1 + \delta}$, under
which the forward rate itself is a \omterm{def:qm:probability-at-speed:martingale}{martingale}; price a swaption under the \omterm{def:qm:girsanov-and-changes-of-numeraire:annuity}{annuity measure} and
compare with a risk-neutral simulation of the whole curve.
\end{tutorial}

\section{Build: moving between numeraires}\label{bld:qm:girsanov-and-changes-of-numeraire:numeraire}

\begin{build}
\bfield{Purpose} Every pricing engine of the miniature firm states its \omterm{def:qm:girsanov-and-changes-of-numeraire:numeraire}{numeraire}, and every
simulated price is tested against it: a price divided by the \omterm{def:qm:girsanov-and-changes-of-numeraire:numeraire}{numeraire} must be a \omterm{def:qm:probability-at-speed:martingale}{martingale}.
\bfield{Interface} \texttt{change\_of\_numeraire\_weights(m\_T, m\_0, n\_T, n\_0)};
\texttt{reweighted\_mean(x, w)}; \texttt{girsanov\_weights(dw, gamma, dt)};
\texttt{deflated\_martingale\_test(prices, numeraire)} returning the largest $|t|$ over dates.
\bfield{Rules} Weights have mean one under the sampling measure (a test, not an assumption); the
integrand of a Girsanov density is adapted; the \omterm{def:qm:probability-at-speed:martingale}{martingale} test uses paths simulated under the
\omterm{def:qm:girsanov-and-changes-of-numeraire:numeraire}{numeraire}'s own measure.
\bfield{Acceptance tests} \texttt{code/firm/numeraire/tests/}: Cameron--Martin weights shift a
Gaussian's mean and not its variance; reweighting risk-neutral short-rate paths reproduces the
forward-measure mean; discounted bond prices pass the \omterm{def:qm:probability-at-speed:martingale}{martingale} test, undiscounted ones fail.
\bfield{Stretch} The spot measure of a discretely rolled account (One Quant Book~6); likelihood-ratio
sensitivities, which are Girsanov weights differentiated in a parameter (One Quant Book~5,
chapter~23).
\end{build}

\begin{omsources}
\item I.~V.~Girsanov, ``On transforming a certain class of stochastic processes by absolutely
continuous substitution of measures'', \emph{Theory of Probability and its Applications} 5, 1960.
\item R.~H.~Cameron and W.~T.~Martin, ``Transformations of Wiener integrals under translations'',
\emph{Annals of Mathematics} 45, 1944.
\item H.~Geman, N.~El Karoui and J.-C.~Rochet, ``Changes of numéraire, changes of probability
measure and option pricing'', \emph{Journal of Applied Probability} 32, 1995.
\item F.~Jamshidian, ``An exact bond option formula'', \emph{Journal of Finance} 44, 1989.
\item W.~Margrabe, ``The value of an option to exchange one asset for another'', \emph{Journal of
Finance} 33, 1978.
\end{omsources}

\section{Exercises}

\begin{exercise}[$\star$]\label{exo:qm:girsanov-and-changes-of-numeraire:1}
Under $d\mathbb Q/d\P = \exp(\gamma W_T - \gamma^2T/2)$ with $\gamma = 0.5$ and $T = 4$, what is the law of $W_T$?
\end{exercise}

\begin{exercise}[$\star$]\label{exo:qm:girsanov-and-changes-of-numeraire:2}
A stock has expected return 9\% and volatility 25\%, and the short rate is 4\%. What is its drift under
the \omterm{def:qm:girsanov-and-changes-of-numeraire:emm}{risk-neutral measure}, and what $\gamma$ does Girsanov use?
\end{exercise}

\begin{exercise}[$\star$]\label{exo:qm:girsanov-and-changes-of-numeraire:3}
A payoff at five years has $\E^{5}[X] = 2.4$. With the chapter's $P(0, 5)$, what is it worth today?
\end{exercise}

\begin{exercise}[$\star\star$]\label{exo:qm:girsanov-and-changes-of-numeraire:4}
Explain why $\E^{\mathbb Q}[r_T] > f(0, T)$ in the chapter's model, and compute the gap at $T = 5$ from
$\sigma^2B(T)^2/2$.
\end{exercise}

\begin{exercise}[$\star\star$]\label{exo:qm:girsanov-and-changes-of-numeraire:5}
Show that the par swap rate is a \omterm{def:qm:probability-at-speed:martingale}{martingale} under the \omterm{def:qm:girsanov-and-changes-of-numeraire:annuity}{annuity measure}, and write the value of a
receiver swaption as an expectation under it.
\end{exercise}

\begin{exercise}[$\star\star$]\label{exo:qm:girsanov-and-changes-of-numeraire:6}
Two stocks follow \omterm{def:qm:stochastic-differential-equations:gbm}{geometric Brownian motions} with volatilities 20\% and correlation 0.5, both at 100,
no dividends. With the second as \omterm{def:qm:girsanov-and-changes-of-numeraire:numeraire}{numeraire}, price the option to receive the first in exchange for
the second in one year.
\end{exercise}

\begin{exercise}[$\star\star\star$]\label{exo:qm:girsanov-and-changes-of-numeraire:7}
Derive the first-passage probability with drift of \cref{prop:qm:brownian-motion:driftpassage}: remove
the drift with Girsanov, apply the reflection principle's joint law of the minimum and the end
point, and integrate. Check it on the stop of chapter~2 with the drift $-\sigma^2/2$.
\end{exercise}

\begin{exercise}[$\star\star\star$]\label{exo:qm:girsanov-and-changes-of-numeraire:8}
\emph{Find the flaw.} ``To price the option we simulate the stock with our research team's
forecast drift of 12\% a year and discount the payoffs at the risk-free rate: the forecast is our
edge.''
\end{exercise}

\section{Problem: Two Measures, One Price}

\begin{problem}[{Weekend problem --- a five-year caplet, two ways}]\label{pb:qm:girsanov-and-changes-of-numeraire:1}
The short rate is $dr = \kappa(\bar r - r)\,dt + \sigma\,dW^{\mathbb Q}$ with $r_0 = 3\%$, $\bar r = 4\%$, $\kappa = 0.5$ and
$\sigma = 1\%$. A caplet fixes at $T_1 = 5$ on the rate for $[5, 5.25]$, pays at 5.25, strike 4.5\%,
notional USD~100 million.

\medskip\noindent\textbf{Part I --- The curve.}
\begin{enumerate}
\item What are $P(0, 5)$ and $P(0, 5.25)$?
\item What is the forward rate for $[5, 5.25]$?
\item What are $\E^{\mathbb Q}[r_5]$ and $f(0, 5)$, and why do they differ?
\item Write the caplet as a put on a zero-coupon bond: strike and number of bonds.
\item What is the bond volatility $\sigma_P$ in Jamshidian's formula?
\end{enumerate}

\medskip\noindent\textbf{Part II --- Three prices.}
\begin{enumerate}[resume]
\item What does the closed form give, in basis points and in dollars?
\item How many normals does a risk-neutral path with weekly steps use, and a forward-measure path?
\item What do 20\,000 risk-neutral paths give, with standard error?
\item And 20\,000 forward-measure paths?
\item Are the two within their standard errors of the closed form?
\end{enumerate}

\medskip\noindent\textbf{Part III --- What the measure change bought.}
\begin{enumerate}[resume]
\item What is the density $d\mathbb Q^{T_1}/d\mathbb Q$ on a path?
\item What does reweighting 100\,000 risk-neutral paths give for $\E^{T_1}[r_5]$?
\item What is the ratio of the two standard errors at equal paths, and why is it close to one?
\item What, then, is the gain at equal accuracy?
\item With four times finer steps under $\mathbb Q$, does the risk-neutral price change?
\end{enumerate}

\medskip\noindent\textbf{Part IV --- Judgement.}
\begin{enumerate}[resume]
\item A cap is a strip of caplets with different payment dates. Which measure do you use?
\item Which measure prices a swaption, and what becomes a \omterm{def:qm:probability-at-speed:martingale}{martingale}?
\item What should a pricing library's acceptance test check about its \omterm{def:qm:girsanov-and-changes-of-numeraire:numeraire}{numeraire}?
\item State the \emph{named result}: the caplet's price and what the \omterm{def:qm:girsanov-and-changes-of-numeraire:forward}{forward measure} bought.
\item In one sentence: what does a \omterm{def:qm:girsanov-and-changes-of-numeraire:changenum}{change of numeraire} change, and what does it not?
\end{enumerate}
\end{problem}

\section{Interview questions}

\begin{interviewq}[$\star$ \iqroles{researcher, bank}]\label{iq:qm:girsanov-and-changes-of-numeraire:1}
What does Girsanov's theorem change about a \omterm{def:qm:brownian-motion:bm}{Brownian motion}, and what can it not change?
\end{interviewq}

\begin{interviewq}[$\star$ \iqroles{bank, trader}]\label{iq:qm:girsanov-and-changes-of-numeraire:2}
Why does a stock drift at the risk-free rate under the \omterm{def:qm:girsanov-and-changes-of-numeraire:emm}{risk-neutral measure}, when nobody expects it to?
\end{interviewq}

\begin{interviewq}[$\star\star$ \iqroles{bank}]\label{iq:qm:girsanov-and-changes-of-numeraire:3}
What is the \omterm{def:qm:girsanov-and-changes-of-numeraire:forward}{forward measure}, and why is the forward rate a \omterm{def:qm:probability-at-speed:martingale}{martingale} under it?
\end{interviewq}

\begin{interviewq}[$\star\star$ \iqroles{bank, researcher}]\label{iq:qm:girsanov-and-changes-of-numeraire:4}
Price the option to exchange asset 2 for asset 1 at $T$. Which \omterm{def:qm:girsanov-and-changes-of-numeraire:numeraire}{numeraire}, and why?
\end{interviewq}

\begin{interviewq}[$\star\star$ \iqroles{developer, bank}]\label{iq:qm:girsanov-and-changes-of-numeraire:5}
How would you test that a Monte Carlo pricing engine uses a consistent measure?
\end{interviewq}

\begin{interviewq}[$\star\star\star$ \iqroles{researcher}]\label{iq:qm:girsanov-and-changes-of-numeraire:6}
Estimate $\P(W_1 > 5)$ by simulation with fewer than 10\,000 draws.
\end{interviewq}
